\section{The alternate collapse and Mayer--Vietoris route}

The nearby-cycle/Leray proof in \cref{sec:integral-leray} gives one verified
route to the topology of $X$.  This section audits the manuscript's independent
cellular-collapse/Mayer--Vietoris route.  It now closes the entire alternate
certificate: the equivariant cusp retraction, its primitive integral cellular
map, equality with analytic nearby-cycle specialization, both geometrically
realized Mayer--Vietoris stages, and the comparison of their obstruction maps
with Leray.

We first correct a conflation in the earlier audit.  For any two-set
decomposition $Y=A\cup B$, our Mayer--Vietoris convention is
\[
\Delta_q=(i_{A*},-i_{B*}):
H_q(A\cap B;\mathbb Z)\longrightarrow
H_q(A;\mathbb Z)\oplus H_q(B;\mathbb Z).
\tag{4.1}
\]
The manuscript does not identify its printed $\alpha_q$ with a one-shot map
for three collars.  It uses two stages.  First,
\[
X^\circ=X_1\cup X_2,
\qquad X_1\cap X_2\simeq F,
\qquad
\alpha_q=(\pi_{1*},-\pi_{2*}).
\tag{4.2}
\]
Second, it attaches the cusp through
\[
X=N'_0\cup X^\circ,
\qquad N'_0\cap X^\circ=M,
\qquad
K_q=\ker\bigl(H_q(M)\to H_q(W)\bigr).
\tag{4.3}
\]
The proof below realizes both stages by actual oriented inclusions and then
also records the equivalent one-shot degree-one collar matrix.

\subsection{The abstract cusp-retraction lemma}

\begin{proposition}[equivariant topological spine]\label{prop:topological-spine}
\Verified\ Let $A$ be a contractible manifold with boundary and let
$A_0\subset\partial A$ be a contractible boundary component.  Let a
group $\Gamma$ act freely and properly discontinuously on $(A,A_0)$.  Assume
$Q=A/\Gamma$ and $H=A_0/\Gamma$ are manifolds, $Q$ is compact,
$H\hookrightarrow Q$ is collared, and the two quotient maps are connected
coverings with deck group $\Gamma$.  If the inclusion identifies the two
deck groups by the identity automorphism of $\Gamma$, then
$H\hookrightarrow Q$ is a strong deformation retract.  The retraction lifts
uniquely from the identity lift to a $\Gamma$-equivariant strong deformation
retraction of $A$ onto $A_0$.
\end{proposition}

\begin{proof}
Because $A$ and $A_0$ are contractible universal covers,
$\pi_n(Q)=\pi_n(H)=0$ for $n\ge2$, while
$\pi_1(Q)\cong\Gamma\cong\pi_1(H)$.  By hypothesis the inclusion induces the
identity on $\Gamma$, so it is a weak homotopy equivalence.  Manifolds have
the homotopy type of CW complexes; Whitehead's theorem
\cite[Theorem~4.5]{Hatcher2002} therefore makes the inclusion a homotopy
equivalence.

The collar makes $(Q,H)$ an NDR pair, hence the inclusion is a cofibration.
We record the standard cofibration argument.  Choose a homotopy inverse
$r_0:Q\to H$.  Since $r_0|_H\simeq\operatorname{id}_H$, the homotopy
extension property modifies $r_0$ through maps $Q\to H$ to a retraction
$r$ satisfying $r|_H=\operatorname{id}_H$.  The maps
$\operatorname{id}_Q$ and $i\circ r$ are homotopic and agree on $H$.
Applying the homotopy extension property once more makes that homotopy
stationary on $H$.  It is therefore a strong deformation retraction.

Let $P:A\to Q$ be the covering.  Lift the deformation homotopy from the
initial lift $\operatorname{id}_A$.  Since the base homotopy fixes $H$ and
the initial lift fixes $A_0$, uniqueness of path lifting makes the lifted
homotopy fix $A_0$.  For $\gamma\in\Gamma$, both
$\widetilde H_t\circ\gamma$ and $\gamma\circ\widetilde H_t$ lift the same
map and agree at $t=0$; uniqueness gives equality for every $t$.  The lift is
therefore $\Gamma$-equivariant and ends in $A_0$.
\end{proof}

\subsection{The candidate-specific equivariant collapse}

We now verify every hypothesis of \cref{prop:topological-spine} in the
candidate.  Retain the toric threefold $Y$, the compact torus
\[
 T_c=(S^1)^3,\qquad T_f=(S^1)^2\times\{1\}\subset T_c,
\]
and the free properly discontinuous action
$\Gamma=\{\Psi_{\bar\lambda}:\bar\lambda\in\bar\Lambda\}$ from
\cref{thm:candidate-cusp-geometry}.  Put
\[
 Y_{\leq\eta}=\{|t|\leq\eta\},\qquad
 Y_{\geq0}=\overline{(\mathbb R_{>0})^3}\subset Y,\qquad
 A=(Y_{\leq\eta})_{\geq0},\qquad A_0=(Y_0)_{\geq0}.
\tag{C.1}
\]
For a torus point write $y=\log|x|/\log|t|$.  The period matrix on a
punctured cusp fibre is
\[
 Z_t=sB_0+C(t),\qquad s=\frac{\log t}{2\pi i},\qquad
 B_t=B_0+\frac{R(t)}{\log|t|},\quad R(t)=-2\pi\operatorname{Im}C(t).
\tag{C.2}
\]
Here $B_0:\bar\Lambda\to\Lambda_{\rm tor}$ is unimodular.

\begin{lemma}[straightening the twist]\label{lem:collapse-straightening}
\Verified\ For all sufficiently small $\eta>0$ there is a homeomorphism
$G:Y_{\leq\eta}\to Y_{\leq\eta}$, fixing $Y_0$ pointwise and preserving
$t$, which conjugates $\Gamma$ to the action $\Gamma^0$ obtained by replacing
$C(t)$ by $C_0=C(0)$.  It commutes with $T_f$ and, on every smooth quotient
fibre, induces a real-linear isomorphism which is the identity on the marked
lattice $\Lambda$.
\end{lemma}

\begin{proof}
Shrink $\eta$ until $B_t$ and
$B_t^0=B_0+R(0)/\log|t|$ are invertible for $0<|t|\leq\eta$.  In logarithmic
coordinates $\xi=(2\pi i)^{-1}\log x$ define
\[
 L_t(\xi)=\xi-(C(t)-C_0)(\operatorname{Im}Z_t)^{-1}
                         \operatorname{Im}\xi.
\tag{C.3}
\]
Every $\xi$ has the unique expression $\xi=Z_ta+b$, with
$a,b\in\mathbb R^2$, and then $L_t(\xi)=Z_t^0a+b$.  Thus $L_t$ is invertible,
fixes $\mathbb R^2$ and its sublattice $\mathbb Z^2$, and sends each deck
translation $Z_t\bar\lambda$ to $Z_t^0\bar\lambda$.  It therefore descends
on $Y_t=\mathbb C^2/\mathbb Z^2$ and conjugates the two $\bar\Lambda$-actions.
In toric coordinates it is
\[
 G(x,t)=\bigl(\exp(-2\pi i(C(t)-C_0)B_t^{-1}y)x,t\bigr),
 \qquad G|_{Y_0}=\operatorname{id}.
\tag{C.4}
\]
It commutes with $T_f$ because $y$ is $T_f$-invariant.

Only continuity at $Y_0$ is nonformal.  Around any $p\in Y_0$, choose one of
the bounded-position subcharts used in the properness proof of
\cref{thm:candidate-cusp-geometry}.  On that subchart $y$ stays bounded,
$B_t^{-1}$ stays bounded, and $C(t)-C_0\to0$.  Hence the torus factor in
(C.4) tends to $1$.  The same argument applies to the inverse, obtained by
interchanging $C(t)$ and $C_0$.  Thus $G$ is a homeomorphism.  Because it
fixes the real lattice and sends the two period columns $Z_t\bar\lambda$ to
$Z_t^0\bar\lambda$, it induces the identity on the marked $H_1=\Lambda$.
\end{proof}

Write $c^0_{\bar\lambda}=u_{\bar\lambda}a_{\bar\lambda}$ for its polar
decomposition, with $u_{\bar\lambda}\in T_f$ and
$a_{\bar\lambda}=\exp(R(0)\bar\lambda)\in(\mathbb R_{>0})^2$.  Let
$\Gamma^+$ be the positive action
\[
 \Psi^+_{\bar\lambda}=(a_{\bar\lambda},1)\Phi_{\bar\lambda}.
\tag{C.5}
\]

\begin{lemma}[positive quotient and honeycomb]\label{lem:positive-collapse}
\Verified\ The pair $(A,A_0)$ in (C.1) consists of contractible topological
manifolds, $A_0$ being a boundary component of $A$.  The action $\Gamma^+$ is
free and properly discontinuous, and
\[
 Q^+=A/\Gamma^+,\qquad H=A_0/\Gamma^+
\tag{C.6}
\]
are respectively a compact $3$-manifold with boundary and one of its boundary
$2$-tori.  The inclusion $H\hookrightarrow Q^+$ is a strong deformation
retract, and its retraction lifts to a $\Gamma^+$-equivariant strong
deformation retraction $P$ of $A$ onto $A_0$.
\end{lemma}

\begin{proof}
Polar multiplication $m:T_c\times Y_{\geq0}\to Y$ is a closed quotient map:
in every smooth unimodular toric chart it is the usual map
$(S^1)^3\times[0,\infty)^3\to\mathbb C^3$, and $T_c$ is compact.  If $q$ is
in the orbit indexed by a face $C$ of the $A_2$ triangulation, its stabilizer
is the connected compact subtorus on
\[
 N_{\widehat C}=\bigoplus_{v\in V(C)}\mathbb Z(v,1).
\]
Unimodularity makes this lattice saturated, and direct intersection with
$\Lambda_{\rm tor}$ gives
\[
 \operatorname{Stab}_{T_c}(q)\cap T_f=
 \begin{cases}
  1,&C\text{ is a vertex},\\
  S^1_{v'-v},&C=\{v,v'\}\text{ is an edge},\\
  T_f,&C\text{ is a triangle}.
 \end{cases}
\tag{C.7}
\]
There are no hidden finite components.  The lattices
$N_{\widehat C}$ and $\Lambda_{\rm tor}$ are saturated, and every face contains
a vertex $(v,1)$, so
$N_{\widehat C}+\Lambda_{\rm tor}=N'$.  Thus the displayed intersections are
the connected subtori attached to their primitive intersection lattices.

In a toric chart the positive set $A$ is
\[
 \{(r_0,r_1,r_2)\in[0,\infty)^3:r_0r_1r_2\leq\eta\}.
\]
Its corner links are discs, so it is a topological $3$-manifold with boundary
$A_0\sqcup\{|t|=\eta\}_{>0}$.  Its interior is
$\mathbb R_y^2\times(0,\eta)$.  Brown's collaring theorem
\cite{Brown1962} makes the inclusion of the interior a homotopy equivalence;
therefore $A$ is contractible.

The nonnegative part of each component $E_v\simeq\mathrm{dP}_6$ is its moment
hexagon.  These hexagons meet along exactly the arcs prescribed by the dual
honeycomb $\mathcal H$ of the $A_2$ triangulation.  Map the two vertex orbits
of $\mathcal H$ to the corresponding torus fixed points, choose homeomorphisms
on representatives of the three edge orbits, and extend the resulting circle
homeomorphism over one hexagon.  Propagating by $B_0\bar\Lambda$ defines a
$\Gamma^+$-equivariant continuous bijection
\[
 \vartheta:\mathbb R^2\xrightarrow{\sim}A_0,\qquad
 \vartheta(z+B_0\bar\lambda)=
 \Psi^+_{\bar\lambda}\vartheta(z).
\tag{C.8}
\]
Local finiteness of the fan makes this bijection proper, hence a homeomorphism.
Thus $A_0$ is contractible.

The action in (C.5) is the already verified cusp action with constant twist
$C^+=(2\pi i)^{-1}R(0)$, since
$e^{2\pi iC^+\bar\lambda}=a_{\bar\lambda}$; hence it is free and properly
discontinuous.  Properness of its quotient over
the disc, proved by the same bounded fundamental-domain argument as
\cref{thm:candidate-cusp-geometry}, makes $Q^+$ compact.  Its two boundary
components are $H\simeq\mathbb R^2/B_0\bar\Lambda$ and
$\mathbb R^2/B_\eta^+\bar\Lambda$, both tori.  Since a free properly
discontinuous action preserves the boundary charts, $Q^+$ is a topological
$3$-manifold with boundary; Brown's collaring theorem makes
$H\hookrightarrow Q^+$ collared.  The covers $A\to Q^+$ and
$A_0\to H$ have the same deck group, and the equivariant inclusion identifies
that group by the identity.  The result is now
\cref{prop:topological-spine}.
\end{proof}

\begin{lemma}[spreading the positive retraction by phases]
\label{lem:phase-spread-collapse}
\Verified\ The lifted retraction $P$ of \cref{lem:positive-collapse} spreads
to a $\Gamma^0$- and $T_c$-equivariant strong deformation retraction of
$Y_{\leq\eta}$ onto $Y_0$.  Conjugating by $G$ gives a $\Gamma$- and
$T_f$-equivariant strong deformation retraction, and hence a strong
deformation retraction
\[
 r_s:N'_0=Y_{\leq\eta}/\Gamma\longrightarrow W=Y_0/\Gamma.
\tag{C.9}
\]
\end{lemma}

\begin{proof}
For $\phi\in T_c$ and $q\in A$ set
\[
 R_s(\phi q)=\phi P_s(q).
\tag{C.10}
\]
This is well defined.  On the open orbit and the outer boundary the
stabilizer is trivial; on $A_0$ the homotopy fixes $q$, so (C.7) causes no
ambiguity.  It is continuous because its pullback along the closed quotient
map $m\times\operatorname{id}_{[0,1]}$ is continuous.

The frozen and positive actions satisfy
\[
 \Psi^0_{\bar\lambda}(\phi q)=
 u_{\bar\lambda}\varphi_{\bar\lambda}(\phi)
 \Psi^+_{\bar\lambda}(q),
\tag{C.11}
\]
where $\varphi_{\bar\lambda}$ is the compact-torus automorphism induced by
$(z,z_3)\mapsto(z+z_3B_0\bar\lambda,z_3)$.  Equation (C.11) and
$\Gamma^+$-equivariance of $P$ prove $\Gamma^0$-equivariance of $R$.
The remaining deformation-retraction identities follow directly from those
of $P$.  Finally \cref{lem:collapse-straightening} conjugates $R$ to the
original action and permits descent to (C.9).
\end{proof}

\begin{proposition}[fixed-fibre cellular collapse]
\label{prop:verified-cellular-collapse}
\Verified\ Fix $t_0$ with $0<|t_0|\leq\eta$ and put $F=f^{-1}(t_0)$.  The
retraction in \cref{lem:phase-spread-collapse} can be chosen so that its end
map $r:N'_0\to W$ restricts to a cellular collapse $c=r|_F$ with the following
description.  There is a splitting
\[
 F\xrightarrow[\bar G]{\sim}F_{t_0}^0\simeq T_B\times T_f,
\tag{C.12}
\]
where $\bar G$ is the identity on the marking $\Lambda$, and the quotient
honeycomb $\Theta\subset T_B$ has two vertices and three edges $e_k$.  Over a
vertex, $c$ collapses all of $T_f$ to a triple point of $W$; over
$e_k^\circ$ it collapses the primitive circle $S^1_{n_k}$ and maps onto the
double curve $\bar C_k$; over $T_B\setminus\Theta$ it is injective onto
$W\setminus D$.
\end{proposition}

\begin{proof}
Let $S\subset A$ be the positive slice $|t|=|t_0|$ and put
\[
 B_\rho^+=B_0+\frac{R(0)}{\log\rho},\qquad
 \ell_\rho=B_0(B_\rho^+)^{-1}.
\]
The map $y:S\to\mathbb R^2$ identifies the positive action with translations
$y\mapsto y+B_{|t_0|}^+\bar\lambda$, while $\ell_{|t_0|}$ conjugates these to
$z\mapsto z+B_0\bar\lambda$.  Define on the frozen fibre
\[
 \widetilde c((\phi,\arg t_0)q)=
 (\phi,\arg t_0)\,
 \vartheta(\ell_{|t_0|}y(q)),\qquad \phi\in T_f, q\in S.
\tag{C.13}
\]
The stabilizer on $S$ is trivial, so the formula is well defined.  Equations
(C.8) and (C.11) make it $\Gamma^0$-equivariant, hence it descends to
$c_\vartheta:F_{t_0}^0\to W$.

Writing
\[
 K_{t_0}=\operatorname{Re}C_0+\frac{\arg t_0}{2\pi}B_0,
\]
the frozen action on $T_f\times\mathbb R_y^2$ is
\[
 \bar\lambda(\phi,y)=
 (\phi e^{2\pi iK_{t_0}\bar\lambda},
  y+B_{|t_0|}^+\bar\lambda).
\]
The change of variables
\[
 \Xi(\phi,y)=
 (\phi e^{-2\pi iK_{t_0}(B_{|t_0|}^+)^{-1}y},y)
\tag{C.14}
\]
conjugates this to the product action and gives (C.12).  Formula (C.7) now
proves the asserted cell description, including primitivity of the collapsed
circle.

It remains to force the end map of an actual strong deformation retraction to
equal (C.13), rather than merely be homotopic to it.  In the chart
$\vartheta:\mathbb R^2\to A_0$ put
\[
 \widehat a=\vartheta^{-1}P_1.
\]
Put $\rho_0=|t_0|$.  If $\rho_0<\eta$, choose a closed positive band
$\{\rho_-\leq |t|\leq\rho_+\}$, disjoint from $A_0$, with $\rho_0$ in its
interior.  If $\rho_0=\eta$, instead use the one-sided band
$\{\rho_-\leq |t|\leq\eta\}$.  Define
$\widehat b=\ell_{|t|}y$ only on the chosen band.  Wherever they are defined,
both maps satisfy
$\widehat a(\Psi^+_{\bar\lambda}q)=\widehat a(q)+B_0\bar\lambda$ and the same
identity for $\widehat b$.  Choose a continuous tent function $\chi(|t|)$,
extended by zero off the band, with $\chi(\rho_0)=1$ and vanishing near every
boundary of the band across which it must be pasted; in the one-sided case no
vanishing is required at the ambient boundary $|t|=\eta$.  With
$w(q)=\chi(|t(q)|)$ define
\[
 K_u(q)=\vartheta\bigl((1-uw(q))\widehat a(q)
                         +uw(q)\widehat b(q)\bigr)
\tag{C.15}
\]
inside the band and $K_u=P_1$ outside.  Because the transition action is by
translations, affine interpolation preserves $\Gamma^+$-equivariance.  The
weight vanishes near every gluing boundary and near $A_0$, so (C.15) pastes
continuously, fixes $A_0$, and takes values in $A_0$.

Concatenate the first half of $P$ with $K_u$.  The result $P'$ still satisfies
all four properties used in \cref{lem:phase-spread-collapse}, and its end map
on $S$ is exactly $\vartheta\ell_{|t_0|}y$.  Spreading $P'$ by phases and
conjugating by $G$ therefore yields a strong deformation retraction whose
restriction to $F$ is
\[
 r|_F=c_\vartheta\circ\bar G.
\]
This proves the proposition, including continuity across all toric strata.
\end{proof}

\subsection{The integral collapse map}

Let $D=\bar C_1\cup\bar C_2\cup\bar C_3$ be the double locus.  Choose
primitive orientations of its three curve classes and of the honeycomb edges
so that the annulus degrees below have the displayed signs.  These orientations
are compared with the clockwise collar convention in
\cref{prop:actual-finite-band-maps,prop:cusp-mv-blocks}.  In
$\Lambda_{\rm tor}=\mathbb Z\langle\widehat w,\widehat\delta\rangle$ put
\[
 n_1=\widehat w,\qquad n_2=\widehat\delta-\widehat w,\qquad
 n_3=-\widehat\delta,\qquad
 \nu_k(v)=\det(n_k,v).
\tag{C.16}
\]

\begin{proposition}[integral homology of the cellular collapse]
\label{prop:verified-collapse-homology}
\Verified\ The map $c_*$ of \cref{prop:verified-cellular-collapse} is
surjective in every degree, and
\[
 0\longrightarrow\operatorname{im}(\textstyle\bigwedge^qM_0-I)
 \longrightarrow\textstyle\bigwedge^q\Lambda
 \xrightarrow{c_*}H_q(W;\mathbb Z)\longrightarrow0
\tag{C.17}
\]
is an exact sequence of free abelian groups.  In the splitting
$\Lambda=\bar\Lambda\oplus\Lambda_{\rm tor}$,
\[
 c_*:\Lambda\to H_1(W)=\bar\Lambda
\]
is projection, and in degree two
\[
 c_*(\widehat\gamma\widehat u)=[\mathbb T_B],\qquad
 c_*(\widehat w\widehat\delta)=0,
\tag{C.18}
\]
\[
 c_*(\beta\otimes v)=
 \sum_{k=1}^3\nu_k(B_0\beta)\nu_k(v)[\bar C_k].
\tag{C.19}
\]
Equivalently, in the ordered source basis
\[
(\widehat\gamma\widehat u,\widehat\gamma\widehat w,
 \widehat\gamma\widehat\delta,\widehat u\widehat w,
 \widehat u\widehat\delta,\widehat w\widehat\delta)
\]
and target basis
$([\mathbb T_B],[\bar C_1],[\bar C_2],[\bar C_3])$,
\[
[c_*]_2=
\begin{pmatrix}
1&0&0&0&0&0\\
0&0&-1&0&0&0\\
0&-1&-1&1&1&0\\
0&0&0&1&0&0
\end{pmatrix}.
\tag{C.20}
\]
In degree three,
\[
 c_*(\widehat\gamma\widehat u\,v)=v
 \quad(v\in\Lambda_{\rm tor}),
 \qquad
 c_*(\beta\,\widehat w\widehat\delta)=0
 \quad(\beta\in\bar\Lambda),
\tag{C.21}
\]
where the first formula fixes the displayed orientation of $H_3(W)$.
\end{proposition}

\begin{proof}
First identify the target integrally.  The normalization
$\nu:\mathrm{dP}_6\to W$ identifies opposite sides of the anticanonical
hexagon.  A plumbing neighbourhood of that hexagon descends to a regular
neighbourhood of $D$.  If $P_0$ is its complement in $\mathrm{dP}_6$, then
$P_0\simeq(\mathbb C^*)^2\simeq T^2$, and excision followed by Lefschetz
duality gives
\[
 H_k(W,D)\simeq H_k(P_0,\partial P_0)
 \simeq H^{4-k}(T^2).
\tag{C.22}
\]
The class $[\mathbb T_B]=c_*[T_B\times\{*\}]$ pairs once with a compact
$T_f$-orbit in $W\setminus D$, so it lifts the generator of $H_2(W,D)$.
Likewise the cycles $c_*(T_B\times S^1_{\widehat w})$ and
$c_*(T_B\times S^1_{\widehat\delta})$ pair unimodularly with the two
coordinate circles in $T_f$.  Thus the boundary maps of the pair vanish and
\[
 H_*(W)=(\mathbb Z,\mathbb Z^2,\mathbb Z^4,\mathbb Z^2,\mathbb Z),
\tag{C.23}
\]
with
$H_2(W)=\mathbb Z\langle[\bar C_1],[\bar C_2],[\bar C_3],
[\mathbb T_B]\rangle$ and $H_3(W)$ generated by the two cycles in (C.21).

The cell description gives projection in degree one and (C.18).  Give $T_B$
the CW structure with one open hexagonal $2$-cell and $1$-skeleton $\Theta$.
A class $\beta\in H_1(T_B)$ has the unique cellular representative
\[
 \sum_{k=1}^3\nu_k(B_0\beta)e_k.
\]
On $e_k\times S^1_v$, the collapse kills $S^1_{n_k}$ and maps the quotient
circle to $\bar C_k$ with degree $\nu_k(v)$.  This proves (C.19), without a
rank inference.  The three choices
$(B_0\beta,v)=(\widehat w,\widehat w),
(\widehat\delta,\widehat\delta),(\widehat w,\widehat\delta)$ give coefficient
vectors
\[
 (0,1,1),\qquad(1,1,0),\qquad(0,1,0),
\]
whose determinant is a unit.  Hence degree two is surjective.  Formula
(C.21) follows directly by crossing the base $2$-cell with a primitive fibre
circle; it proves degree-three surjectivity and records the orientations.
 In degree four the top-dimensional open stratum
$(T_B\setminus\Theta)\times T_f$ maps homeomorphically to the open
$4$-stratum of $W$, so the map is multiplication by a unit after choosing the
target orientation.  Degree zero is immediate.

Let $h$ be geometric monodromy around the cusp.  Fibre transport in $N'_0$
is a homotopy from the inclusion of $F$ to its composite with $h$.  Composing
with $r$ yields $c\simeq c\circ h$, and therefore
\[
 \operatorname{im}(\textstyle\bigwedge^qM_0-I)\subseteq\ker c_*.
\tag{C.24}
\]
The inclusion is an equality integrally, not only rationally.  In the ordered
basis $(\widehat\gamma,\widehat u,\widehat w,\widehat\delta)$, direct exterior
multiplication gives the primitive image lattices
\[
\begin{array}{c|c}
q&\operatorname{im}(\bigwedge^qM_0-I)\\ \hline
0,4&0\\
1&\langle\widehat w,\widehat\delta\rangle\\
2&\langle\widehat w\widehat\delta,
          \widehat\gamma\widehat w+\widehat u\widehat\delta\rangle\\
3&\langle\widehat\gamma\widehat w\widehat\delta,
          \widehat u\widehat w\widehat\delta\rangle.
\end{array}
\tag{C.25}
\]
Each displayed subgroup is a direct summand.  Its rank is
$\binom4q-b_q(W)$ by (C.23), while surjectivity gives the same rank for
$\ker c_*$.  The quotient $H_q(W)$ is free, so $\ker c_*$ is saturated.
Two saturated sublattices of the same rank and rational span coincide, proving
(C.17).
\end{proof}

\subsection{Comparison with analytic nearby cycles}

\begin{theorem}[collapse equals nearby-cycle specialization]
\label{thm:collapse-equals-specialization}
\Verified\ With the radial marking of the punctured cusp family, for every
$0\leq q\leq4$ one has an equality of integral homomorphisms
\[
 \operatorname{sp}^q=i_F^*r^*=c^*:
 H^q(W;\mathbb Z)\longrightarrow H^q(F;\mathbb Z).
\tag{C.26}
\]
Consequently the cellular calculation gives
\[
 \operatorname{im}c^*=(\textstyle\bigwedge^qV)^{T_0}
\]
as a saturated lattice.  In particular, the degree-three orientation in the
cellular computation is the degree-three orientation of intrinsic
specialization; there is no independent sign or index.
\end{theorem}

\begin{proof}
Let $N_\rho=f^{-1}(|t|\leq\rho)\subset N'_0$ and let
$i_{W,\rho}:W\hookrightarrow N_\rho$ be the central-fibre inclusion.  The
closed tubes are cofinal among neighbourhoods of $W$: if an open set
$\mathcal O\subset N'_0$ contains $W$, properness makes
$f(N'_0\setminus\mathcal O)$ closed and disjoint from $0$, so
$N_\rho\subset\mathcal O$ for all sufficiently small $\rho$.  For each such
$\rho$, choose $t_\rho$ on the radial ray through the fixed point $t_0$ whose
fibre is $F$, with $0<|t_\rho|<\rho$, and write
$i_{t_\rho}:F_{t_\rho}\hookrightarrow N_\rho$.  Identify
$H^q(F_{t_\rho};\mathbb Z)$ with $H^q(F;\mathbb Z)$ by parallel transport
along that ray.  The neighborhood realization of the unit
$\mathbb Z_W\to\psi_f\mathbb Z$, together with proper base change, identifies
specialization with
\[
 H^q(W;\mathbb Z)
 \xleftarrow[\sim]{\ \rho_0\ }
 \varinjlim_{\rho\downarrow0}H^q(N_\rho;\mathbb Z)
 \xrightarrow{\ \rho_{\rm near}\ }H^q(F;\mathbb Z),
\tag{C.27}
\]
where $\rho_0$ is induced by $i_{W,\rho}^*$ and $\rho_{\rm near}$ is induced
by the transported maps $i_{t_\rho}^*$.  These restriction maps are compatible
under shrinking $\rho$.  This is the same analytic integral
construction used in
\cref{thm:p0-integral-specialization}; see
\cite[Definition~4.2.1 and Propositions~4.2.2, 4.2.11]{Dimca2004}.

Applying the construction of \cref{lem:phase-spread-collapse} at radius
$\rho$ shows that $i_{W,\rho}^*$ is an isomorphism; no compatibility of those
radius-dependent deformation homotopies is needed.  Let
$j_\rho:N_\rho\hookrightarrow N'_0$ be inclusion.  The single retraction
$r:N'_0\to W$ satisfies
\[
 i_{W,\rho}^*j_\rho^*r^*
 =(r\circ j_\rho\circ i_{W,\rho})^*
 =\operatorname{id}_{H^q(W;\mathbb Z)}.
\]
Hence $j_\rho^*r^*$ is the unique inverse of $i_{W,\rho}^*$.  These inverses
are compatible under shrinking because they are restrictions of the single
map $r^*:H^q(W;\mathbb Z)\to H^q(N'_0;\mathbb Z)$.  Consequently
$\rho_0^{-1}(\alpha)$ in (C.27) is represented, for every sufficiently small
$\rho$, by $j_\rho^*r^*\alpha$.

The restrictions of one class on $N'_0$ to the fibres along the punctured
radial ray form a section of the Gauss--Manin local system.  Therefore, after
the parallel-transport identifications used to define $\rho_{\rm near}$,
$i_{t_\rho}^*j_\rho^*r^*\alpha$ equals $i_F^*r^*\alpha$.  It follows that the
composite in
(C.27) is $i_F^*r^*$.  Since $c=r|_F$, functoriality gives
$i_F^*r^*=(r\circ i_F)^*=c^*$.

Changing the lift of the radial marking composes the target with monodromy.
By (C.24), $c^*$ lands in the monodromy-invariant lattice, so the equality is
independent of that harmless choice.  Finally
\cref{thm:p0-integral-specialization} identifies the common image with the
full saturated invariant lattice.
\end{proof}

\begin{test}[resolved cusp-collapse falsifiers]
\Verified\ The audit tested the failure modes previously attached to C1--C3:
continuity at every toric stratum, finite stabilizer components, equivariance,
primitive edge circles, degree-two and degree-three orientations, torsion and
saturation, and functorial equality with the nearby-cycle unit.  The proofs
above exclude each one.  This status concerns the cusp collapse only; it does
not infer the global Mayer--Vietoris matrices.
\end{test}

\begin{auditrecord}[V1: based van Kampen inclusions]
\label{record:vankampen}
\Verified\ The actual overlap coordinate maps, with the zero section as
base locus and clockwise meridians, show that the
finite fillings impose $\widehat\rho_j^{m_j}=t_{v_j}$ with the same sign,
the cusp filling kills $\Lambda_{\rm tor}$ and its toric meridian, and the
tautological gluing has zero translation drift.  Thus the abstract terminal
presentation is the actual $\pi_1(X)$.  The deck generators, common clockwise
convention, zero-section disc, normal closure, boundary inclusion
$\Lambda\to\Lambda/\Lambda_{\rm tor}$,
$\widetilde a_0\mapsto1$, and absence of extra whisker drift are proved in
\cref{prop:geometric-pi1}.
\end{auditrecord}

\subsection{The actual finite-band Mayer--Vietoris maps}

Fix the ordered orientation
\[
H_1(F)=\Lambda=
\mathbb Z\langle\widehat\gamma,\widehat u,
                    \widehat w,\widehat\delta\rangle
\]
dual to $(\gamma,u,w,\delta)$, use the complex orientations of
$F,S_1,S_2$, and orient every base meridian clockwise.  Put
\[
 \eta_1=2u+w+3\delta,
 \qquad \eta_2=u+w+2\delta,
 \qquad q=uw+6\gamma\delta.
\tag{MV.1}
\]

\begin{proposition}[actual finite-band maps in primitive bases]
\label{prop:actual-finite-band-maps}
\Verified\ In the decomposition $X^\circ=X_1\cup X_2$, the maps denoted
$\alpha_r$ in the manuscript are the actual Mayer--Vietoris maps
\[
 \alpha_r=(\pi_{1*},-\pi_{2*}):
 H_r(F;\mathbb Z)\longrightarrow
 H_r(S_1;\mathbb Z)\oplus H_r(S_2;\mathbb Z).
\tag{MV.2}
\]
There is no conjugacy, unrecorded index, or independent orientation sign in
(MV.2).

Choose primitive bases
\[
 H_1(S_1)=\mathbb Z\langle g_1,h_1\rangle,
 \quad [\widehat\gamma]_1=3g_1,\quad
 h_1=\pi_{1*}\widehat w,
\]
\[
 H_1(S_2)=\mathbb Z\langle g_2,h_2\rangle,
 \quad [\widehat\gamma]_2=-4g_2,\quad
 h_2=\pi_{2*}\widehat w.
\]
Then, in the source basis
$(\widehat\gamma,\widehat u,\widehat w,\widehat\delta)$ and the target
basis $(g_1,h_1,g_2,h_2)$,
\[
 [\alpha_1]=
 \begin{pmatrix}
 3&0&0&0\\
 0&2&1&3\\
 4&0&0&0\\
 0&-1&-1&-2
 \end{pmatrix}.
\tag{MV.3}
\]
It has Smith form $\operatorname{diag}(1,1,1,0)$ and
\[
 \ker\alpha_1=\mathbb Z\nu_0,
 \qquad
 \nu_0=-\widehat u-\widehat w+\widehat\delta.
\tag{MV.4}
\]

Define cohomology bases $(a_j,b_j)$ by
\[
 \pi_1^*a_1=\gamma\eta_1,
 \qquad \pi_1^*b_1=q-2\gamma\eta_1,
\]
\[
 \pi_2^*a_2=2\gamma\eta_2,
 \qquad \pi_2^*b_2=q-3\gamma\eta_2,
\tag{MV.5}
\]
and let $(a_j^\sharp,b_j^\sharp)$ be the Kronecker-dual homology bases.
In the ordered source basis
\[
(\widehat\gamma\widehat u,
 \widehat\gamma\widehat w,
 \widehat\gamma\widehat\delta,
 \widehat u\widehat w,
 \widehat u\widehat\delta,
 \widehat w\widehat\delta)
\]
and target basis
$(a_1^\sharp,b_1^\sharp,a_2^\sharp,b_2^\sharp)$,
\[
 [\alpha_2]=
 \begin{pmatrix}
  2& 1& 3& 0&0&0\\
 -4&-2& 0& 1&0&0\\
 -2&-2&-4& 0&0&0\\
  3& 3& 0&-1&0&0
 \end{pmatrix}.
\tag{MV.6}
\]
This matrix also has Smith form $\operatorname{diag}(1,1,1,0)$; its
primitive left kernel is
\[
 \Phi=(4,2,3,2).
\tag{MV.7}
\]
The class $\mathcal A\in H^2(X^\circ;\mathbb Z)$ represented by $\Phi$
restricts to $2q$, because
\[
 2q=4\pi_1^*a_1+2\pi_1^*b_1
    =3\pi_2^*a_2+2\pi_2^*b_2.
\tag{MV.8}
\]

If the degree-three bases are Poincar\'e dual to the degree-one cohomology
bases just fixed, the projection formula makes $[\alpha_3]$ the same matrix
as (MV.3).  Finally
\[
 [\alpha_0]=\binom{1}{-1},
 \qquad
 [\alpha_4]=\binom{3}{-4}.
\tag{MV.9}
\]
Every image in (MV.3), (MV.6), and (MV.9) is saturated.
\end{proposition}

\begin{proof}
Choose closed discs $D'_1,D'_2$ whose intersection is a contractible band
disjoint from the three special points.  Over the annular shell between
$D'_j$ and the smaller filling disc, $f$ is a proper submersion.  An
Ehresmann connection lifts a base deformation retraction that is stationary
near the smaller disc; adjoining the identity there gives a deformation
retraction $X_j\to N_j$.  The verified local radial retraction of
\cref{lem:finite-filling-stalk-map} keeps the marked flat coordinate $x$
fixed and sends its endpoint to the orbit $\pi_j(x)\in S_j$.  Since the band
is contractible, choose its marking by the same radial paths.  The composite
of the actual band inclusion with $X_j\to N_j\to S_j$ is therefore literally
$\pi_j$.  The standard chain-level Mayer--Vietoris injection is
$z\mapsto(i_{1*}z,-i_{2*}z)$
\cite[\S2.2]{Hatcher2002}, proving (MV.2) with its sign.

The affine fundamental-group presentations of the finite fibres and the
exact finite pullback lattices of \cref{prop:finite-indices} give the stated
primitive degree-one bases and (MV.3).  Its $3\times3$ minors include $3$
and $4$, so their gcd is one; direct elimination gives (MV.4).
The two bases in (MV.5) are unimodular changes of
\[
 \langle\gamma\eta_1,q\rangle,
 \qquad
 \langle2\gamma\eta_2,\gamma\eta_2+q\rangle,
\]
the exact integral images of $\pi_j^*$ in degree two.  Evaluating them on the
six source tori gives (MV.6).  Its rank-three minors include $2$ and $-3$,
so the image is saturated; multiplication verifies (MV.7) and (MV.8).

For degree three, if $x\in H_3(F)$ and $\theta\in H^1(S_j)$, then
\[
 \langle\theta,\pi_{j*}x\rangle
 =\langle\pi_j^*\theta,x\rangle.
\]
Thus in the Poincar\'e-dual bases the degree-three matrix equals the
degree-one matrix.  In degree four, the holomorphic quotient coverings have
positive degrees $3$ and $4$; the sole minus sign is the Mayer--Vietoris
minus sign.  This proves (MV.9) and completes the integral certificate.
\end{proof}

\subsection{The cusp attachment and its primitive blocks}

\begin{proposition}[actual cusp kernels and global blocks]
\label{prop:cusp-mv-blocks}
\Verified\ For the second decomposition $X=N'_0\cup X^\circ$, the map
$H_r(M)\to H_r(W)$ is surjective and
\[
 K_1=\mathbb Z[\widetilde a_0],
 \qquad
 K_2=\mathbb Z\langle
       \xi_{\widehat w},\xi_{\widehat\delta}\rangle,
\tag{MV.10}
\]
where
\[
 \xi_\nu=a_*([S^1_\nu]\times[\widetilde a_0])
\]
uses the oriented fibrewise translation action and the clockwise toric
meridian.  Both displayed subgroups are primitive.

Let $\zeta$ be the generator of $H_1(X^\circ)$ normalized by
$\Phi_1(\zeta)=1$.  In the source basis
$(\widehat\gamma,\widehat u,\widetilde a_0)$ of $H_1(M)$ and target basis
$(\bar\gamma,\bar u,\zeta)$ of
$H_1(W)\oplus H_1(X^\circ)$, the actual map
$(c_*,-j_*)$ is
\[
 D_1=
 \begin{pmatrix}
 1&0&0\\
 0&1&0\\
 -12&0&-p
 \end{pmatrix},
 \qquad \det D_1=-p.
\tag{MV.11}
\]

Orient the connecting map by
\[
 \partial(j_*\xi_\nu)=(I-A_2)\nu=(I-A_1^{-1})\nu
 \quad(\nu\in\Lambda_{\rm tor}).
\tag{MV.12}
\]
Let $e$ be the primitive generator of
$C_2:=\operatorname{coker}\alpha_2$ with
$\langle\mathcal A,e\rangle=1$, and put
$z=j_*\xi_{\widehat w}$.  Then $(e,z)$ is a basis of
$H_2(X^\circ)$ and
\[
 j_*\xi_{\widehat\delta}=-p\,e,
 \qquad
 [j_*|_{K_2}]_{(\xi_{\widehat\delta},\xi_{\widehat w})}^{(e,z)}
 =\begin{pmatrix}-p&0\\0&1\end{pmatrix}.
\tag{MV.13}
\]
After splitting the surjections onto $H_r(W)$, the full degree-one and
degree-two Mayer--Vietoris maps are integrally equivalent to
\[
 I_2\oplus[p],
 \qquad
 I_4\oplus
 \begin{pmatrix}-p&0\\0&1\end{pmatrix},
\tag{MV.14}
\]
respectively.  They are unimodular for the candidate $p=-1$.
\end{proposition}

\begin{proof}
The composite $H_r(F)\to H_r(M)\to H_r(W)$ is the verified collapse $c_*$,
which is surjective by \cref{prop:verified-collapse-homology}; hence the
second arrow is surjective.  Under the Wang projection, the base class
$[\widetilde a_0]$ is primitive.  It bounds the explicit zero-section disc
in the toric filling, proving the first formula in (MV.10).

For $K_2$, use the toric chart $t=z_1z_2z_3$.  The literal solid tori
\[
 \{|z_1|\leq\rho,\ |z_2|=r,\ z_3=c\}
 \quad\hbox{and}\quad
 \{|z_2|\leq\rho,\ |z_1|=r,\ z_3=c\}
\]
have oriented boundaries $\xi_{\widehat w}$ and
$\xi_{\widehat\delta}$.  Under
\[
 H_2(M)\twoheadrightarrow
 \ker(M_0-I|\Lambda)=\Lambda_{\rm tor},
\]
these classes map to $\widehat w$ and $\widehat\delta$.  Their span is
therefore primitive; the Wang ranks and the verified homology of $W$ show
that it is all of $K_2$.

The actual meridian calculation gives
\[
 j_*\widehat\gamma=12\zeta,
 \qquad j_*\widehat u=0,
 \qquad j_*[\widetilde a_0]=p\zeta,
\]
while $c_*$ is projection onto
$\mathbb Z\langle\bar\gamma,\bar u\rangle$.  This is (MV.11).
Because $A_1A_2=M_0^{-1}$ acts trivially on $\Lambda_{\rm tor}$, the two
expressions in (MV.12) agree.  Since
$A_2\widehat w=-\widehat u+\widehat\delta$ and
$A_2\widehat\delta=\widehat\delta$,
\[
 \partial(j_*\xi_{\widehat w})=-\nu_0,
 \qquad \partial(j_*\xi_{\widehat\delta})=0.
\tag{MV.15}
\]
Thus $z$ is a primitive lift, up to the fixed unit sign, of the generator of
$\ker\alpha_1$.

It remains to determine the coefficient of $e$.  Normalize contraction by
\[
 \iota_v(\theta_1\wedge\cdots\wedge\theta_k)
 =\sum_r(-1)^{r-1}\theta_r(v)
   \theta_1\wedge\cdots\widehat{\theta_r}\cdots\wedge\theta_k.
\]
The globally free $S^1_{\widehat\delta}$-action and the sweep/slant identity
give
\[
 \langle\mathcal A,j_*\xi_{\widehat\delta}\rangle
 =\langle\iota_{\widehat\delta}\mathcal A,
          j_*[\widetilde a_0]\rangle,
 \qquad
 \iota_{\widehat\delta}(2q)=-12\gamma.
\tag{MV.16}
\]
Let $\Theta\in H^1(X^\circ)$ be the primitive class restricting to
$12\gamma$.  Restriction to $F$ is injective, so (MV.16) implies
$\iota_{\widehat\delta}\mathcal A=-\Theta$.  The based meridian relation
gives
\[
 \Theta(j_*[\widetilde a_0])=p.
\]
Hence the left side of (MV.16) is $-p$, which proves (MV.13).  Splitting the
surjective cusp blocks gives (MV.14).
\end{proof}

\begin{remark}[one-shot $H_1$ check]
The preceding two-stage calculation is the manuscript's argument.  There is
also a useful one-shot check.  Let $U$ be the torus bundle over the
pair-of-pants core and
$N=N'_0\sqcup N'_1\sqcup N'_2$.  In the target basis
\[
(c,\rho_1,\rho_2;
 \bar\gamma,\bar u;
 g_1,h_1;g_2,h_2)
\]
and domain basis
\[
(\gamma_0,u_0,a_0;
 \gamma_1,w_1,\rho_1;
 \gamma_2,w_2,\rho_2),
\]
the actual map $(i_{U*},-i_{N*})$ for the candidate is
\[
\Delta_1=
\begin{pmatrix}
1&0&\ell_0&1&0&0&1&0&0\\
0&0&-1&0&0&1&0&0&0\\
0&0&-1&0&0&0&0&0&1\\
-1&0&0&0&0&0&0&0&0\\
0&-1&0&0&0&0&0&0&0\\
0&0&0&-3&0&-1&0&0&0\\
0&0&0&0&-1&0&0&0&0\\
0&0&0&0&0&0&4&0&-1\\
0&0&0&0&0&0&0&-1&0
\end{pmatrix}.
\tag{MV.17}
\]
For example, $a_0$ maps to
$\ell_0c-\rho_1-\rho_2$ on the pair-of-pants side and to zero on the cusp
side; $\gamma_1$ maps to $c$ and $3g_1$, while $\gamma_2$ maps to $c$ and
$-4g_2$.  Thus every column is an actual primitive collar cycle.  Exact
elimination gives
\[
 \det\Delta_1=1-12\ell_0=-p
\]
for $(\ell_1,\ell_2)=(1,-1)$.  At $\ell_0=0$, this determinant is one.
\end{remark}

\subsection{Exact comparison with the Leray transgressions}

\begin{theorem}[Mayer--Vietoris/Leray comparison]
\label{thm:mv-leray-comparison}
\Verified\ With the preceding primitive bases and clockwise meridians, the
actual Mayer--Vietoris obstruction maps commute with the intrinsic Leray
transgressions in degrees one, two, and three.  If
$\omega_{\rm MV}\in H^2(B;\mathbb Z)$ is the generator fixed below, then
\[
 d_2^{0,1}(12\gamma)=-p\,\omega_{\rm MV},
\tag{L.1}
\]
\[
 d_2^{0,2}(2q)=-p\,\omega_{\rm MV}[\delta],
 \qquad
 d_2^{0,3}(2\gamma uw)
   =-p\,\omega_{\rm MV}[\gamma\delta].
\tag{L.2}
\]
For the candidate $p=-1$, all three maps are multiplication by $+1$ in
these bases.
\end{theorem}

\begin{proof}
Let $V$ be the pair-of-pants neighbourhood in the base and let
$D_i$ be the three special discs.  Order the annuli as
$(D_1\cap V,D_2\cap V,D_0\cap V)$ and orient their meridians clockwise.
Define
\[
 \omega_{\rm MV}=\partial_B[(1,0,0)].
\]
The chain-level Mayer--Vietoris sequence on the base gives
\[
 \partial_B(n_1,n_2,n_0)
 =(n_1+n_2+n_0)\omega_{\rm MV},
\tag{L.3}
\]
because the image of $H^1(V)$ is the sum-zero sublattice.

Choose a lift $\Theta_V$ of $12\gamma$ over the pair of pants and write
$b_i=\Theta_V(\widehat\rho_i)$.  The finite local lifts have meridian values
$4\ell_1$ and $3\ell_2$, while the cusp lift vanishes on
$\widetilde a_0$.  For the general affine cusp regluing used in this
comparison, let $t_\mu$ denote translation by the class
$\mu\in\Lambda/\Lambda_{\rm tor}$ and set
$\ell_0=\gamma(\mu)$.  Thus
$\Theta_V(t_\mu)=12\ell_0$.  The tautological gluing of the constructed
candidate has $\mu=0$ and $\ell_0=0$; in that case the following relation
has no vertical drift.  In general,
\[
 \widetilde a_0=t_\mu\widehat\rho_2^{-1}
                         \widehat\rho_1^{-1},
\]
the local-minus-pair-of-pants discrepancy is
\[
 \bigl(4\ell_1-b_1,\ 3\ell_2-b_2,\
 -(12\ell_0-b_1-b_2)\bigr).
\tag{L.4}
\]
Its signed sum is $-p$.  The canonical truncation triangle identifies this
very cocycle, functorially on the induced total-space cover, with Leray
$d_2$; this is the construction proved in
\cref{lem:first-transgression}.  Equations (L.3)--(L.4) prove (L.1).

This is more than equality of the terminal integer.  The transpose of the
oriented collar map applied to the tuple of local lifts has meridian values
\[
 (b_1-4\ell_1,\ b_2-3\ell_2,\
 12\ell_0-b_1-b_2),
\]
the negative of (L.4).  Consequently the square
\[
\begin{CD}
\mathbb Z\langle12\gamma\rangle
 @>{d_2^{0,1}}>>
\mathbb Z\langle\omega_{\rm MV}\rangle\\
@V{\cong}VV @VV{\cong}V\\
H^1(X^\circ;\mathbb Z)
@>{(j_*|_{K_1})^\vee}>>
K_1^\vee
\end{CD}
\tag{L.5}
\]
commutes when the left vertical map sends $12\gamma$ to $\Theta$ and the
right sends $\omega_{\rm MV}$ to $-\widetilde a_0^\vee$.  The lower map is
multiplication by $p$, while the upper map is multiplication by $-p$.

For degree two, quotient the primitive lift in (MV.15):
\[
 \overline K_2=K_2/\mathbb Z\xi_{\widehat w},
 \qquad C_2=\operatorname{coker}\alpha_2.
\]
Equation (MV.13) gives an actual reduced map
\[
 \overline j:\overline K_2\longrightarrow C_2,
 \qquad [\xi_{\widehat\delta}]\longmapsto -p\,e.
\tag{L.6}
\]
Under the primitive identifications
\[
 2q\longmapsto\Phi\in C_2^\vee,
 \qquad
 \omega_{\rm MV}[\delta]\longmapsto
 \xi_{\widehat\delta}^{\,\vee}\in\overline K_2^\vee,
\]
the square
\[
\begin{CD}
\mathbb Z\langle2q\rangle
@>{d_2^{0,2}}>>
\mathbb Z\langle\omega_{\rm MV}[\delta]\rangle\\
@V{\cong}VV @VV{\cong}V\\
C_2^\vee @>{\overline j^{\,\vee}}>> \overline K_2^\vee
\end{CD}
\tag{L.7}
\]
commutes, and both horizontal maps are multiplication by $-p$.  This proves
the first formula in (L.2).

The final formula has no new sign choice.  Write
\[
 d_2(2q)=r\,\omega_{\rm MV}[\delta],
 \qquad
 d_2(2\gamma uw)=s\,\omega_{\rm MV}[\gamma\delta].
\]
Using
\[
 (12\gamma)(2q)=12(2\gamma uw),
 \qquad (2q)(2\gamma uw)=0,
\]
the derivation rule and the already fixed degree-one sign give
\[
 s=2(-p)-r.
\]
The identities
\[
 \delta\wedge2\gamma uw=-2\operatorname{vol},
 \qquad
 2q\wedge\gamma\delta=2\operatorname{vol}
\]
applied to the second product give $r=s$.  Hence $r=s=-p$, proving the
second formula in (L.2).

Changing every meridian orientation reverses $p$ and
$\omega_{\rm MV}$ simultaneously.  Reversing the
$S^1_{\widehat\delta}$ orientation reverses both
$\xi_{\widehat\delta}$ and its dual coefficient generator.  Thus all
remaining convention changes are common units and cannot change an index.
\end{proof}

\begin{test}[resolved global topological falsifiers]
\Verified\ The actual finite-band maps are (MV.2), their degree-one and
degree-two images are saturated, the multiple-fibre top class is
$(3,-4)$, the cusp kernels are literal primitive collar cycles, the reduced
cusp blocks are (MV.11) and (MV.13), and the comparison is the pair of
commuting squares (L.5) and (L.7), not a comparison of absolute determinants.
The former MV1 and L1 failure modes are therefore excluded for the candidate.
\end{test}

\subsection{Complete recognition theorem}

\begin{lemma}[cohomological closure in dimension six]
\label{lem:cohomological-closure}
\Verified\ Let $X$ be a smooth, closed, connected, oriented six-manifold.
If
\[
\pi_1(X)=0,
\qquad H^2(X;\mathbb Z)=H^3(X;\mathbb Z)=0,
\tag{4.4}
\]
then $H_k(X;\mathbb Z)$ is $\mathbb Z$ for $k=0,6$ and zero for
$1\le k\le5$.
\end{lemma}

\begin{proof}
The Hurewicz map in degree one identifies
$H_1(X;\mathbb Z)=\pi_1(X)^{\mathrm{ab}}=0$.  The universal coefficient
sequence
\[
0\longrightarrow\operatorname{Ext}(H_1(X),\mathbb Z)
\longrightarrow H^2(X;\mathbb Z)
\longrightarrow\operatorname{Hom}(H_2(X),\mathbb Z)
\longrightarrow0
\]
shows that $\operatorname{Hom}(H_2(X),\mathbb Z)=0$, so the finitely
generated group $H_2(X)$ is torsion.  The next universal coefficient
sequence contains an injection
\[
\operatorname{Ext}(H_2(X),\mathbb Z)
\hookrightarrow H^3(X;\mathbb Z)=0.
\]
For a finitely generated torsion group
$A=\bigoplus_i\mathbb Z/n_i$, one has
$\operatorname{Ext}(A,\mathbb Z)=\bigoplus_i\mathbb Z/n_i$; hence
$H_2(X)=0$.  Poincar\'e duality gives
\[
H_3(X)\simeq H^3(X)=0,\quad
H_4(X)\simeq H^2(X)=0,\quad
H_5(X)\simeq H^1(X)=0.
\]
Finally connectedness and orientation give $H_0(X)=H_6(X)=\mathbb Z$.
\end{proof}

\begin{theorem}[sphere recognition from the exact inputs]
\label{thm:sphere-recognition}
\Verified\ Let $X$ be a smooth, closed, connected six-manifold.  If
\[
\pi_1(X)=0,\qquad
H_k(X;\mathbb Z)\cong
\begin{cases}
\mathbb Z,&k=0,6,\\
0,&1\le k\le5,
\end{cases}
\tag{4.5}
\]
then $X$ is diffeomorphic to the standard $S^6$.
\end{theorem}

\begin{proof}
A closed smooth manifold has the homotopy type of a finite CW complex.
We first prove successively that $\pi_i(X)=0$ for $1\le i\le5$.
The case $i=1$ is assumed.  Suppose $2\le n\le5$ and
$\pi_i(X)=0$ for $i<n$.  Then $X$ is $(n-1)$-connected, so the Hurewicz
homomorphism
\[
h_n:\pi_n(X)\longrightarrow H_n(X;\mathbb Z)
\]
is an isomorphism by the Hurewicz theorem
\cite[Theorem~4.32]{Hatcher2002}.  The target vanishes by (4.5), hence
$\pi_n(X)=0$.  Induction proves the assertion.

Now $X$ is $5$-connected.  Hurewicz in degree six gives
$\pi_6(X)\cong H_6(X;\mathbb Z)\cong\mathbb Z$.  Choose a map
$g:S^6\to X$ whose homotopy class maps to a chosen orientation generator.
It has degree one.  Hence $g_*$ is an isomorphism in degrees zero and six;
in every other degree both integral homology groups vanish.  Thus $g$ is an
integral homology equivalence.  Both spaces are simply connected CW
complexes, so the homological Whitehead theorem
\cite[Corollary~4.33]{Hatcher2002} makes $g$ a homotopy equivalence.
Therefore $X$ is a smooth homotopy $6$-sphere.

Smale's generalized Poincar\'e theorem in dimensions greater than four
\cite[Theorem~A]{Smale1961} identifies the underlying topological manifold
with $S^6$.  The smooth endpoint is stated directly by Smale:
every homotopy $6$-sphere is diffeomorphic to $S^6$
\cite[Corollary~1.2]{Smale1962}.  Equivalently, smooth homotopy
$6$-spheres up to oriented $h$-cobordism form $\Theta_6$;
Kervaire--Milnor's opening table gives $|\Theta_6|=1$, and their following
remark identifies $h$-cobordism with diffeomorphism in dimension six
\cite[opening table and following Remark, pp.~504--505]{KervaireMilnor1963}.
Hence there is no exotic smooth structure in dimension six, and $X$ is
diffeomorphic to the standard sphere.
\end{proof}

\begin{corollary}[abstract application theorem]\label{cor:abstract-s6}
\Verified\ Let $X$ be a smooth, closed, connected, oriented six-manifold.
Suppose V1 identifies the actual based van Kampen quotient
with $G(\ell_0,\ell_1,\ell_2)$ of
\cref{prop:terminal-group}, and suppose an independently constructed
MV1/L1 comparison proves
\[
H^2(X;\mathbb Z)\cong H^3(X;\mathbb Z)
\cong\mathbb Z/|p|\mathbb Z,
\qquad p=12\ell_0-4\ell_1-3\ell_2.
\tag{4.6}
\]
If $|p|=1$, then $X\cong_{\mathrm{diff}}S^6$.
\end{corollary}

\begin{proof}
For $|p|=1$, 
\cref{prop:terminal-group} and the assumed geometric identification give
$\pi_1(X)=0$, while (4.6) gives $H^2(X)=H^3(X)=0$.
\Cref{lem:cohomological-closure} converts these inputs into the homology
condition (4.5), and
\cref{thm:sphere-recognition} gives the claimed diffeomorphism.
\end{proof}

\begin{corollary}[independent Mayer--Vietoris certificate for the candidate]
\label{cor:candidate-mv-sphere}
\Verified\ For the tautological candidate
$(\ell_0,\ell_1,\ell_2)=(0,1,-1)$, the two-stage geometric
Mayer--Vietoris calculation gives
\[
 H_k(X;\mathbb Z)\cong
 \begin{cases}
 \mathbb Z,&k=0,6,\\
 0,&1\leq k\leq5.
 \end{cases}
\]
Together with the independently verified based van Kampen calculation, this
proves $X\cong_{\mathrm{diff}}S^6$ without using the Leray computation of the
same groups.
\end{corollary}

\begin{proof}
Here $p=-1$.  Thus (MV.11) and the reduced block (MV.13) are unimodular.
The two Mayer--Vietoris sequences give $H_1(X)=H_2(X)=0$ integrally, with no
torsion hidden by a rank calculation.  The verified collapse homology
(C.23) gives $e(W)=1-2+4-2+1=2$.  Away from $W$, the map to the punctured
base is stratified by smooth four-tori and the two bielliptic fibres, all of
Euler characteristic zero.  Additivity therefore gives $e(X)=e(W)=2$.

Poincar\'e duality and universal coefficients give $H_5(X)=H_4(X)=0$.
Moreover
\[
 H_3(X)\cong H^3(X)\cong\operatorname{Hom}(H_3(X),\mathbb Z),
\]
because $H_2(X)=0$; hence $H_3(X)$ is free.  The equality $e(X)=2$ forces
its rank to be zero.  The based van Kampen result
\cref{prop:geometric-pi1} gives $\pi_1(X)=1$, and
\cref{thm:sphere-recognition} completes the proof.
\end{proof}

Thus C1--C3, V1, MV1, and L1 are all verified for the candidate.  The
nearby-cycle/Leray route and the collapse/Mayer--Vietoris route are now two
separate integral certificates for the same sphere-recognition input.
